\begin{frame}{Problem, który rozwiązuje praca}
  \begin{columns}[T,onlytextwidth]
    \begin{column}{0.54\textwidth}
      \begin{itemize}
        \item Krótko opisz kontekst i praktyczny problem.
        \item Wskaż lukę w istniejących rozwiązaniach.
        \item Wyjaśnij, dlaczego problem jest istotny.
      \end{itemize}
      \vspace{3mm}
      \begin{block}{Główna teza}
        Jedno zdanie, które słuchacze powinni zapamiętać z tego slajdu.
      \end{block}
    \end{column}
    \begin{column}{0.42\textwidth}
      \imageplaceholder[0.38\textheight]{Miejsce na rysunek problemu, fotografię stanowiska albo schemat systemu}
    \end{column}
  \end{columns}
\end{frame}

\begin{frame}{Cel i zakres pracy}
  \begin{columns}[T,onlytextwidth]
    \begin{column}{0.48\textwidth}
      \textbf{Cel główny}
      \begin{itemize}
        \item Zdefiniuj mierzalny rezultat pracy.
      \end{itemize}
      \vspace{3mm}
      \textbf{Zakres}
      \begin{itemize}
        \item analiza wymagań,
        \item projekt i implementacja,
        \item badania oraz ocena wyników.
      \end{itemize}
    \end{column}
    \begin{column}{0.48\textwidth}
      \textbf{Kryteria sukcesu}
      \begin{enumerate}
        \item Kryterium jakościowe lub ilościowe.
        \item Warunek brzegowy albo ograniczenie.
        \item Sposób weryfikacji rezultatu.
      \end{enumerate}
    \end{column}
  \end{columns}
\end{frame}
